\documentclass{tutordesk}

\studio[NE]{Northbridge English Studio}
\student{Mina Park}
\lesson{8}
\lessondate{2 Nov 2026}
\doctype{Worksheet}
\exam{IELTS \& PTE Academic \textperiodcentered{} Mixed skills}
\title{Mixed Practice: Reading, Listening and Writing}
\subtitle{Six short tasks in exam format, about 55 minutes}
\level{IELTS 6.0 to 7.0 / PTE 58 to 65}
\timing{55 min}
\focus{Scanning, spelling, cohesion}

\begin{document}
\maketitle

% ---------------------------------------------------------------- 1
\begin{exercise}[tag=IELTS Reading]{True, False or Not Given}
\instructions{Read the passage, then decide whether statements 1--5 agree with the information in it.}

\begin{passage}[Cooling the City][Practice text written for this worksheet. All people named are fictional.]
\plabel{A}On a still summer afternoon, the centre of a large city can be several degrees
warmer than the farmland around it. This effect, known as the urban heat island, is caused
less by the sun itself than by what cities are made of. Asphalt, brick and concrete absorb
heat during the day and release it slowly at night, so streets that feel bearable at noon
can remain uncomfortable well after midnight.

\plabel{B}For many years the standard response was simply more air conditioning. Yet
air-conditioning units pump warm air back into the street, and the electricity they use
often comes from power stations that add to the problem. Urban climatologist Dr Ines Varga
describes this as \enquote{cooling the inside by heating the outside} and argues that it
can only ever be a short-term measure.

\plabel{C}Planners have therefore turned to passive solutions. The simplest is paint: roofs
coated in pale, reflective material stay markedly cooler than dark ones. Street trees are
more expensive to maintain, but they provide shade as well as cooling through evaporation
from their leaves. According to landscape architect Tom\'as Reyes, a single mature tree can
do more for a pavement than any new technology, \enquote{provided someone remembers to
water it in its first three summers}.

\plabel{D}Green roofs, covered in soil and low-growing plants, have attracted the most
attention. They insulate the buildings beneath them and absorb rainwater that would
otherwise overload drains. However, they are heavy, and older buildings may need structural
work before one can be installed. Building engineer Priya Nair points out that the cost of
this work is rarely included in the optimistic figures quoted by suppliers.

\plabel{E}None of these measures is a complete answer on its own. Researchers increasingly
speak of a \enquote{cooling network}: parks, tree-lined streets and reflective surfaces
linked so that cool air can move through a neighbourhood. Whether cities have the patience
to build such networks, which take decades rather than years to mature, remains an open
question.
\end{passage}

\begin{truefalse}[TFNG]
\begin{twocol}
  \statement{Cities are warmer than the countryside mainly because they receive more sunlight.}{FALSE}
  \why{Para A: ``less by the sun itself than by what cities are made of''.}
  \statement{Air conditioning can make the street outside a building warmer.}{TRUE}
  \why{Para B: units ``pump warm air back into the street''.}
  \statement{Reflective roof paint is now required by law in several cities.}{NOT GIVEN}
  \statement{Street trees cost more to look after than reflective roofs.}{TRUE}
  \why{Para C: trees are ``more expensive to maintain'' than paint.}
  \statement{Residents generally prefer green roofs to street trees.}{NOT GIVEN}
  \why{Para D says green roofs attract attention, not who prefers them.}
\end{twocol}
\end{truefalse}
\end{exercise}

% ---------------------------------------------------------------- 2
\begin{exercise}[tag=IELTS Reading]{Matching features}
\instructions{Look at the people in the box and the statements 6--8. Match each statement with the correct person, A, B or C.}
\begin{matching}[People]
  \moption{Dr Ines Varga, urban climatologist}
  \moption{Tom\'as Reyes, landscape architect}
  \moption{Priya Nair, building engineer}
  \mitem{A low-cost measure may fail without basic care in its early years.}{B}
  \mitem{Suppliers tend to understate the full cost of a measure.}{C}
  \mitem{One common solution moves a problem rather than solving it.}{A}
\end{matching}
\end{exercise}

% ---------------------------------------------------------------- 3
\begin{exercise}[tag=IELTS Listening]{Note completion}
\instructions{Your teacher will read a short talk once. Complete the notes. Write ONE WORD AND/OR A NUMBER for each answer.}

\begin{teachernote}[Audio script, read once at natural speed]
Hello, and welcome to Riverside Community Garden. I'm Josie Marsh, that's M-A-R-S-H, and
I look after the volunteers. Please don't wait at the main gate: we meet at the greenhouse,
right next to the car park. Sessions run every Saturday from half past nine. We used to
finish at twelve, but we now carry on until twelve thirty. You don't need tools, we have
plenty, but do bring gloves, and in summer a hat, because there's very little shade. For your
first few weeks you'll mostly be weeding and watering the vegetable beds; pruning comes
later, after some training. Teenagers are very welcome, but anyone under sixteen needs to
come with an adult.
\end{teachernote}

\noindent{\sffamily\bfseries Riverside Community Garden: volunteer induction}
\begin{gapfill}
\begin{twocol}
\begin{itemize}[leftmargin=1.1em,itemsep=4pt]
  \item Meet at the \gap{greenhouse} next to the car park
  \item Saturdays, 9.30 to \gap[18mm]{12.30}
  \item Bring gloves and a \gap{hat}
  \item New volunteers: weeding and \gap{watering}
  \item Under \gap[14mm]{16}: come with an adult
  \item Coordinator: Josie \gap{Marsh}
\end{itemize}
\end{twocol}
\end{gapfill}
\end{exercise}

% ---------------------------------------------------------------- 4
\begin{exercise}[tag=PTE Reading]{Fill in the blanks}
\instructions{Choose the word from the box that best completes each gap. There are three words you do not need.}
\wordbank{implications, reinforce, abundant, random, contradicts, consolidate, evidence, reduces}
\begin{gapfill}
Sleep researchers have long known that memory is not simply recorded during the day and
stored at night. Instead, the sleeping brain appears to \gap{consolidate} recent experiences,
strengthening some connections and allowing others to fade. This process is not
\gap{random}: information linked to a reward, or that the learner expected to need later,
is more likely to be retained. The \gap{implications} for students are considerable.
Revising late into the night may feel productive, but if it \gap{reduces} the total amount
of sleep, it can undermine the very learning it was meant to \gap{reinforce}.
\end{gapfill}
\end{exercise}

% ---------------------------------------------------------------- 5
\begin{exercise}[tag=PTE Reading]{Re-order paragraphs}
\instructions{The sentences below are in a random order. Write the letters in the correct order.}
\begin{reorder}
  \para{4}{By marking each death on a street map, he showed that cases clustered around a single public water pump.}
  \para{2}{At the time, most doctors believed the disease was spread through foul air.}
  \para{5}{His map is now regarded as one of the founding documents of modern epidemiology.}
  \para{1}{In 1854, a cholera outbreak killed hundreds of people within a few streets of central London.}
  \para{3}{The physician John Snow, however, suspected that contaminated water was to blame.}
\end{reorder}
\why{D opens with time and place; ``At the time'' (B) needs a date before it; ``however'' (E) contrasts with B; ``his map'' (C) needs the map in A. PTE scores each correct adjacent pair.}
\end{exercise}

% ---------------------------------------------------------------- 6
\begin{exercise}[tag=PTE Listening]{Write from dictation}
\instructions{Your teacher will read each sentence once. Write the sentence exactly as you hear it.}
\begin{teachernote}
Read each sentence once at natural speed and pause for about 15 seconds before the next.
\end{teachernote}
\writeline{The library will be closed for maintenance during the first week of term.}
\writeline{Students are expected to submit their assignments through the online portal.}
\writeline{Most of the research was conducted in partnership with local schools.}
\end{exercise}

% ---------------------------------------------------------------- 7
\begin{exercise}[tag=IELTS Writing Task 2]{Plan before you write}
\instructions{Spend about eight minutes planning an answer to the question below. Use one line per move.}

\noindent\emph{Some people think that young people should have to do unpaid community work in
their free time. To what extent do you agree or disagree?}

\subsection{Position}
\answerlines{2}
\subsection{Body paragraph 1: Point, Explain, Example}
\answerlines{3}
\subsection{Body paragraph 2: Point, Explain, Example}
\answerlines{3}
\subsection{Conclusion: link back to the question}
\answerlines{2}
\begin{solution}[Plan]
\emph{Position:} partly agree: community work is valuable, but making it compulsory weakens it.\par
\emph{Body 1:} P: volunteering builds skills school does not teach. Ex: teamwork with adults, responsibility to real people. Eg: a teenager running a food-bank shift learns scheduling and patience.\par
\emph{Body 2:} P: compulsion turns service into a box-ticking exercise. Ex: motivation drops when choice disappears. Eg: required hours often become the minimum, done without interest.\par
\emph{Conclusion:} schools should reward and promote community work rather than require it.
\end{solution}
\end{exercise}

\printanswers
\end{document}
